% CPSY 1291 -- Recitation handout 10
% Compile:  latexmk -lualatex handout-10-sequences-recurrence-dynamics.tex
\documentclass[11pt]{article}
\usepackage{cpsy1291-handout}

\begin{document}

\handouthead{10}{Sequences, recurrence, and dynamics}
  {Unrolling, vanishing gradients, fixed points, and a Hopfield network in fifteen lines}
  {Week 10 \ $\cdot$\ hold this after Lecture 15; Assignment 5 is due Tue 11/24}

\section*{What this session is for}
\addcontentsline{toc}{section}{What this session is for}

Two lectures, one object. Lecture 14 treats a recurrent network as a machine
that processes sequences; Lecture 15 treats the same equations as a
\emph{dynamical system} with fixed points and basins of attraction. They are the
same object viewed from two directions, and the second view is what lets you
reverse-engineer a trained network instead of merely evaluating it.

This session is the code for both, plus the two things that reliably go wrong:
shapes, and gradients that vanish.

\begin{keybox}
\ask{The one idea.} A recurrent network is a \emph{map} $\boldsymbol{h} \mapsto
f(\boldsymbol{h}, \boldsymbol{x})$ applied over and over. Everything else
follows from that: training difficulties come from differentiating a long
composition of that map, and memory comes from the map having states it does not
leave.
\end{keybox}

\section{Shapes, before anything else}

\begin{center}
\begin{tabular}{@{}lll@{}}
\toprule
Object & Shape & Note \\
\midrule
a batch of sequences & $(N, T, D)$ & with \texttt{batch\_first=True} \\
the same, PyTorch default & $(T, N, D)$ & the default is \emph{not} batch-first \\
hidden state & $(\text{layers}, N, H)$ & note: batch is in the middle, always \\
\bottomrule
\end{tabular}
\end{center}

\begin{keybox}
\ask{Always pass \texttt{batch\_first=True}} to \texttt{nn.RNN},
\texttt{nn.LSTM}, and \texttt{nn.GRU}, and note that the hidden state ignores it
--- $h$ is $(\text{layers}, N, H)$ either way. That inconsistency is the single
most common source of silent shape bugs in recurrent code. It is silent because
$(T, N, D)$ and $(N, T, D)$ have the same number of elements, so nothing errors
when $T$ and $N$ happen to be compatible; you simply train on transposed data.
\end{keybox}

\section{A recurrent cell, by hand}

\[
\boldsymbol{h}_t = \tanh(\boldsymbol{W}_{hh}\boldsymbol{h}_{t-1} +
\boldsymbol{W}_{xh}\boldsymbol{x}_t + \boldsymbol{b})
\]

\begin{lstlisting}
class RNNCell(nn.Module):
    def __init__(self, d_in, d_hid):
        super().__init__()
        self.Wxh = nn.Linear(d_in, d_hid)
        self.Whh = nn.Linear(d_hid, d_hid, bias=False)

    def forward(self, x, h):
        return torch.tanh(self.Wxh(x) + self.Whh(h))

h = torch.zeros(N, d_hid)
for t in range(T):
    h = cell(X[:, t], h)          # the SAME weights at every step
\end{lstlisting}

Two things to notice in that loop. The weights do not change with $t$ --- that is
\textbf{weight sharing across time}, the temporal analogue of a convolution's
weight sharing across space, and it is what lets one network handle sequences of
any length. And the loop builds a graph $T$ layers deep: \textbf{unrolling}. An
RNN run for 50 steps is a 50-layer feedforward network whose layers happen to be
identical.

\subsection{The library version}

\begin{lstlisting}
rnn = nn.LSTM(input_size=D, hidden_size=H, num_layers=1, batch_first=True)
out, (h_n, c_n) = rnn(X)          # out: (N, T, H)   h_n: (1, N, H)
\end{lstlisting}

\texttt{out} holds the hidden state at \emph{every} step; \texttt{h\_n} holds only
the last one. Which you want depends on the task: classifying a whole sequence
uses \texttt{h\_n} (or a pooled \texttt{out}); labeling every step uses
\texttt{out}.

\section{Why gradients vanish, and how to see it}

Differentiating a loss at step $T$ with respect to a weight at step $1$ requires
the chain rule through every intervening step, which produces a \emph{product} of
$T$ Jacobians. Products of many numbers below 1 go to zero; products of numbers
above 1 explode. That is the whole phenomenon.

\begin{lstlisting}
loss.backward()
for t, h in enumerate(hidden_states):
    print(t, h.grad.norm().item())      # requires h.retain_grad() beforehand
\end{lstlisting}

Plot those norms against $t$ on a log axis. A clean exponential decay backwards
in time is the vanishing gradient made visible, and it is worth seeing once ---
it converts a piece of received wisdom into an observation.

\begin{center}
\begin{tabular}{@{}ll@{}}
\toprule
Problem & What you do about it \\
\midrule
exploding & \texttt{clip\_grad\_norm\_(model.parameters(), 1.0)} --- reliable, one line \\
vanishing & gating: LSTM or GRU. Clipping does not help \\
\bottomrule
\end{tabular}
\end{center}

\textbf{Gating} works because the cell state passes through an \emph{addition}
rather than a repeated multiplication, so the gradient has a path back through
time that is not a long product. That is the mechanism; the three gates are the
bookkeeping around it.

\subsection{Teacher forcing}

When training a model to generate a sequence, you can feed back either the
model's own output or the true previous token. Feeding the truth is
\textbf{teacher forcing}: it trains far faster and more stably, and it creates a
mismatch --- at generation time the model sees its own outputs, including its own
mistakes, which it never saw during training. Errors then compound. This is worth
knowing before Lecture 17, where the same trick trains every language model.

\section{The same equations as a dynamical system}

Set the input to a constant and the update becomes a map from
$\boldsymbol{h}$ to $\boldsymbol{h}$. Iterate it and ask where it goes.

\subsection{Finding fixed points numerically}

A \textbf{fixed point} satisfies $\boldsymbol{h}^* = f(\boldsymbol{h}^*)$. Find
them by minimizing $q(\boldsymbol{h}) = \tfrac12\|f(\boldsymbol{h}) -
\boldsymbol{h}\|^2$ from many random starting points:

\begin{lstlisting}
h = torch.randn(200, H, requires_grad=True)       # 200 starting points at once
opt = torch.optim.Adam([h], lr=0.01)
for _ in range(2000):
    opt.zero_grad()
    q = 0.5 * ((cell(x_const, h) - h) ** 2).sum(1)
    q.sum().backward()
    opt.step()
found = h[q.detach() < 1e-8]                      # keep the ones that converged
\end{lstlisting}

This is the standard method for reverse-engineering a trained RNN, and it is
short because autograd does the work. Starting from many points matters: a
network can have several fixed points, and each one has its own basin.

\subsection{Stable or unstable?}

Linearize at a fixed point --- take the Jacobian $\boldsymbol{J} = \partial
f/\partial\boldsymbol{h}$ evaluated there --- and look at its eigenvalues.

\begin{center}
\begin{tabular}{@{}ll@{}}
\toprule
All $|\lambda| < 1$ & \textbf{stable}: nearby states fall in. A memory \\
Some $|\lambda| > 1$ & \textbf{unstable} in those directions: a saddle \\
One $|\lambda| \approx 1$ & a \textbf{line attractor}: the state drifts freely along it \\
\bottomrule
\end{tabular}
\end{center}

The third row is the interesting one for neuroscience: a direction that neither
decays nor grows is a way to hold a continuous quantity, which is how a network
integrates evidence or holds a value in working memory.

\begin{lstlisting}
J = torch.autograd.functional.jacobian(lambda h: cell(x_const, h), h_star)
np.abs(np.linalg.eigvals(J.numpy()))
\end{lstlisting}

\section{A Hopfield network in fifteen lines}

The clearest case of a network whose memories \emph{are} its attractors.

\begin{lstlisting}
def store(patterns):                       # patterns: (P, n) of +/-1
    W = patterns.T @ patterns / len(patterns)      # Hebbian outer product
    np.fill_diagonal(W, 0)                         # no self-connections
    return W

def recall(W, s, steps=20, rng=np.random.default_rng(0)):
    s = s.copy()
    for _ in range(steps):
        for i in rng.permutation(len(s)):          # asynchronous update
            s[i] = 1 if W[i] @ s >= 0 else -1
    return s

def energy(W, s):
    return -0.5 * s @ W @ s
\end{lstlisting}

Three points to check by running it. Recall \textbf{monotonically decreases the
energy}, which is why it converges rather than cycling --- and asynchronous
updating is required for that guarantee. The \textbf{capacity} is about
$0.14n$ patterns for $n$ units; store more and recall returns blends. And the
\textbf{negation} of every stored pattern is also a fixed point, which falls
straight out of the symmetry of the update rule.

\section{Exercises}

\ask{1.} You pass $(N, T, D)$ data to \texttt{nn.LSTM} without
\texttt{batch\_first=True}. Nothing errors. What is the model learning?

\ask{2.} Your RNN trains fine on sequences of length 10 and not at all on
sequences of length 200. Name the mechanism, and the architectural fix.

\ask{3.} Gradient clipping at 1.0 does not help a network whose loss is flat
from epoch 1. Why not?

\ask{4.} You find 40 candidate fixed points from 200 random starts, but 160 of
them have $q > 10^{-3}$. What are those 160, and should you report them?

\ask{5.} A trained RNN has a fixed point whose Jacobian has one eigenvalue at
$0.999$ and the rest below $0.6$. What is the network doing?

\ask{6.} A Hopfield network with 100 units is given 30 patterns to store. What
happens on recall, and why?

\subsection*{Answers}

\ask{1.} It reads your batch dimension as time and your time dimension as batch:
each ``sequence'' is one time step across $N$ different examples. It trains, it
produces a loss, and it is learning nothing about temporal structure. Nothing
errors because both layouts contain the same number of elements.

\ask{2.} Vanishing gradients: the gradient from step 200 back to step 1 is a
product of 200 Jacobians. Fix: gating --- an LSTM or GRU --- because the cell state
passes through an addition rather than a repeated multiplication.

\ask{3.} Clipping bounds gradients from above; it does nothing to gradients that
are too small. A flat loss from the start is vanishing gradients, a dead
nonlinearity, or a learning rate far too small --- none of which clipping
addresses.

\ask{4.} They are starting points that did not converge --- the optimizer was
still descending, or it was heading toward a region with no fixed point. Discard
them, and report both numbers: ``40 fixed points found from 200 random
initializations''. Reporting only the 40 hides how thoroughly you searched.

\ask{5.} It has a \textbf{line attractor}: one direction along which the state
neither decays nor grows, and every other direction contracting onto it. That is
how a network holds a continuous quantity --- integrating evidence, or
maintaining a value in working memory.

\ask{6.} Capacity is about $0.14 \times 100 = 14$ patterns, so 30 is roughly
double it. Recall converges to spurious states --- blends and mixtures of stored
patterns --- rather than to the patterns themselves, and the failure gets worse
the more you store.

\section*{Cheat sheet}
\addcontentsline{toc}{section}{Cheat sheet}

\begin{center}
\begin{tabular}{@{}p{0.44\textwidth}p{0.5\textwidth}@{}}
\toprule
\textbf{You want} & \textbf{You write} \\
\midrule
sequences shaped how you think     & \texttt{nn.LSTM(..., batch\_first=True)} \\
every step's hidden state          & \texttt{out, \_ = rnn(X)} $\rightarrow$ \texttt{(N, T, H)} \\
only the last                      & \texttt{\_, (h\_n, \_) = rnn(X)} $\rightarrow$ \texttt{(1, N, H)} \\
to stop exploding gradients        & \texttt{clip\_grad\_norm\_(m.parameters(), 1.0)} \\
to stop vanishing gradients        & use gating; clipping will not help \\
gradient norms through time        & \texttt{h.retain\_grad()}, then \texttt{h.grad.norm()} \\
fixed points                       & minimize $\tfrac12\|f(h)-h\|^2$ from many starts \\
stability at a fixed point         & eigenvalues of the Jacobian; $|\lambda|<1$ is stable \\
a Hopfield weight matrix           & \texttt{P.T @ P / len(P)}, zero the diagonal \\
Hopfield capacity                  & about $0.14n$ patterns \\
\bottomrule
\end{tabular}
\end{center}

\end{document}
